\section{Instalação}

O Docker Desktop está disponível para Mac, Windows e Linux. No Linux, o Docker provê pacotes \texttt{.deb} e \texttt{.rpm} para as seguintes distribuições e arquiteturas: Ubuntu, Debian, Red Hat e Fedora. Existe uma versão experimental para distribuições baseadas em Arch, porém o Docker não testou nem verificou a instalação nessas distribuições.

\subsection{Requisitos do sistema}

Para executar o Docker Desktop no Linux, são necessários os seguintes requisitos:

\begin{itemize}
    \item Kernel 64-bit e suporte a virtualização pela CPU
    \item Suporte a KVM para virtualização
    \item QEMU 5.2 ou superior (recomenda-se a versão mais recente)
    \item Sistema de inicialização \texttt{systemd}
    \item 4 GB de RAM
\end{itemize}

Vale ressaltar que o Docker Desktop no Linux roda uma máquina virtual.

\subsection{Instalação no Ubuntu}

Para instalar o Docker Desktop no Ubuntu, basta baixar o pacote mais recente e executar os seguintes comandos:

\begin{verbatim}
sudo apt-get update
sudo apt install ./docker-desktop-amd64.deb
\end{verbatim}
